\chapter{Pregnancy Test (hCG)}
\section*{What Is This Test?}
A pregnancy test detects human chorionic gonadotropin (hCG), a glycoprotein hormone produced by the syncytiotrophoblast cells of the placenta after implantation (approximately 6--12 days after fertilization) \cite{wilcox1999time}. hCG consists of an alpha subunit (identical to the alpha subunits of LH, FSH, and TSH) and a unique beta subunit that confers biologic and immunologic specificity. The primary function of hCG is to maintain the corpus luteum and its progesterone production during the first 6--8 weeks of gestation until the placenta takes over progesterone synthesis. hCG also stimulates fetal testicular testosterone production and maternal thyroid (TSH cross-reactivity at high levels can cause gestational hyperthyroidism). Pregnancy tests detect the beta subunit of hCG in urine (qualitative, point-of-care [POC] or home pregnancy test) or serum (quantitative, measured by immunoassay) \cite{chard1992pregnancy}. Urine pregnancy tests typically detect hCG $\geq$25 mIU/mL. Serum quantitative hCG is detectable as early as 8--11 days after conception, doubling every 48--72 hours during the first 8--10 weeks, reaching a peak of 50,000--100,000 mIU/mL at 10--12 weeks, declining to a plateau of 10,000--30,000 mIU/mL for the remainder of gestation. hCG measurement is used for: confirmation of pregnancy, evaluation of suspected ectopic pregnancy (serial hCG that fails to double in 48 hours or rises suboptimally suggests ectopic or failed intrauterine pregnancy), diagnosis and monitoring of gestational trophoblastic disease (hydatidiform mole, choriocarcinoma --- hCG levels are extremely elevated, often >100,000 mIU/mL), first-trimester screening for trisomy 21 (elevated free beta-hCG) and trisomy 18 (decreased free beta-hCG), and tumor marker for testicular germ cell tumors (seminoma and non-seminoma --- hCG is elevated in seminoma [15--20\%], choriocarcinoma [always], and embryonal carcinoma).
\section*{Alternative Names}
hCG Test; Beta-hCG; Qualitative hCG; Quantitative hCG; Pregnancy Urine Test; Pregnancy Serum Test
\section*{Principle}
Urine pregnancy test (POC or home): A lateral flow immunochromatographic sandwich immunoassay using anti-beta-hCG monoclonal antibodies. Urine applied to the test strip migrates by capillary action; if hCG is present, it binds to an anti-beta-hCG antibody conjugated to colored particles (colloidal gold or latex), and the hCG-antibody-particle complex is captured by a second anti-beta-hCG antibody immobilized at the test line, producing a colored line. The control line confirms test validity. Serum quantitative hCG: Two-site sandwich chemiluminescent immunoassay (CLIA) or electrochemiluminescent immunoassay (ECLIA) on automated immunoassay platforms, detecting intact hCG and free beta-hCG \cite{cole1997immunoassay}. Quantitative results reported in mIU/mL.
\section*{History}
The first pregnancy test to gain clinical acceptance was the Aschheim--Zondek test, described in 1928, which detected hCG by injecting the patient's urine into immature female mice and observing ovarian changes. In the 1930s, the Friedman test used rabbits, and by the 1960s the frog and toad (Xenopus) test was widely used in clinical laboratories. The first immunologic pregnancy tests, based on hemagglutination inhibition, appeared in the 1960s. The development of monoclonal antibodies specific for the beta subunit of hCG in the 1970s enabled sensitive and specific immunoassays, and the first home pregnancy test was marketed in the United States in 1977. Modern qualitative urine tests use lateral flow immunochromatographic technology, and quantitative serum assays on automated platforms now detect hCG within days of implantation.
\section*{Why This Test Is Ordered}
\begin{itemize}
  \item Confirmation of pregnancy in a woman of reproductive age with amenorrhea, nausea, breast tenderness, or suspected pregnancy
  \item Evaluation of first-trimester bleeding or pelvic pain, where a positive or negative test helps distinguish threatened abortion, ectopic pregnancy, and failed intrauterine pregnancy
  \item Detection and monitoring of ectopic pregnancy: serial quantitative serum hCG every 48 hours to assess doubling time \cite{barnhart2009ectopic}
  \item Diagnosis and monitoring of gestational trophoblastic disease (hydatidiform mole, choriocarcinoma), which produces markedly elevated hCG
  \item First-trimester combined screening for trisomy 21 (Down syndrome) and trisomy 18, where free beta-hCG is measured with PAPP-A and nuchal translucency
  \item Tumor marker for testicular germ cell tumors and, rarely, other hCG-secreting neoplasms
  \item Preoperative or pre-imaging testing to avoid operating on an undiagnosed pregnancy
\end{itemize}
\section*{Is This a Primary or Secondary Test}
A urine pregnancy test is a primary screening test for pregnancy, being rapid, inexpensive, and highly accurate when used after the missed period. Quantitative serum hCG is a secondary, more sensitive test used to confirm and monitor the course of early pregnancy, to evaluate suspected ectopic pregnancy or trophoblastic disease, and to interpret ultrasound findings.
\section*{How This Test Is Performed}
For a home or clinic urine test, the patient voids into a clean container or holds the test stick in the urine stream, and the result is read within the time specified (typically 3--5 minutes) after urine migrates across the immunochromatographic strip. For a quantitative serum test, blood is drawn by standard venipuncture into a plain (red-top) or serum separator tube, and the serum is analyzed on an automated two-site immunoassay platform (chemiluminescent or electrochemiluminescent) that reports hCG in mIU/mL. First morning urine is preferred for qualitative testing because it is most concentrated, although modern urine tests are sensitive enough for random samples once hCG is $\geq$25 mIU/mL.
\section*{What Preparation Is Needed}
No special preparation is required. For urine tests, patients should avoid excessive fluid intake immediately before testing, because dilute urine can cause a false-negative result, and first morning void is preferred. Some medications and fertility treatments containing hCG (such as those used for ovulation induction) can cause a false-positive result, and the patient should inform the clinician about any hCG-containing therapy.
\section*{Contraindications and Precautions}
There are no contraindications to urine or serum hCG testing. Precautions include interpreting the result in clinical context: a negative test in very early pregnancy may reflect hCG below the assay threshold, and a positive test must be correlated with the menstrual history, pelvic ultrasound, and serial hCG values. In suspected ectopic pregnancy, a single positive test has limited value, and the discriminatory zone (hCG $\geq$1,500--2,000 mIU/mL) guides ultrasound interpretation \cite{acog2018tubal}. Some medications and heterophile antibodies can cause false-positive results; if the result is discordant with the clinical picture, it should be confirmed by a different assay or a repeat test.
\section*{Risks}
Minimal risk. The only risks are those of venipuncture for serum testing (brief pain, bruising, and rarely hematoma or infection). Urine testing is entirely noninvasive.
\section*{What Happens After the Test}
A positive qualitative test is usually followed by a confirmatory quantitative serum hCG, especially when ectopic pregnancy or abnormal pregnancy is suspected. If the pregnancy is desired, prenatal care, a dating ultrasound, and serial monitoring as indicated follow. If the hCG does not rise appropriately, a repeat measurement, transvaginal ultrasound, and gynecologic consultation are arranged.
\section*{Understanding Results}
\subsection*{Negative (Not Pregnant)}
Urine: hCG <25 mIU/mL. Serum: hCG <5 mIU/mL. A negative pregnancy test in a woman with amenorrhea, pelvic pain, and risk factors for ectopic pregnancy requires repeat testing in 48 hours, pelvic ultrasound, and consideration of ectopic pregnancy. Early pregnancy may be below the detection limit.
\subsection*{Positive (Pregnant)}
Urine: hCG $\geq$25 mIU/mL. Serum: hCG $\geq$5--25 mIU/mL (depending on assay threshold). Normal early intrauterine pregnancy: hCG doubles every 48--72 hours (minimum 53\% increase in 48 hours). Ectopic pregnancy: hCG rises suboptimally (<53\% increase in 48 hours), plateaus, or rises and falls. Failed/inevitable miscarriage: hCG falls. Gestational trophoblastic disease (hydatidiform mole): markedly elevated hCG (>100,000--200,000 mIU/mL) with the uterus larger than dates and hyperemesis gravidarum.
\subsection*{False-Positive hCG}
Heterophile antibodies and HAMA (human anti-mouse antibodies) cross-linking the capture and detection antibodies in the absence of hCG, producing false-positive results. Confirmed by: negative urine hCG (heterophile antibodies are too large to filter into urine), serial dilutions showing non-linear results, or hCG measurement by a different assay platform. Menopause: pituitary hCG production (physiologic low-level hCG <10--15 mIU/mL) in postmenopausal women from pituitary gonadotroph cells. Quiescent gestational trophoblastic disease.
\section*{Normal Results}
Non-pregnant: Serum hCG <5 mIU/mL (not pregnant). Urine hCG: negative. Pregnancy (quantitative serum hCG): 3 weeks LMP (approximately 1 week post-conception): 5--50 mIU/mL. 4 weeks: 10--500 mIU/mL. 5 weeks: 200--7,000 mIU/mL. 6 weeks: 200--32,000 mIU/mL. Peak at 10--12 weeks: 50,000--100,000+ mIU/mL. The absolute hCG level is less important than the doubling time in early pregnancy. A single hCG value cannot distinguish intrauterine from ectopic pregnancy.
\section*{Follow-up Tests}
Serial quantitative serum hCG (every 48 hours), transvaginal pelvic ultrasound (visualization of gestational sac at hCG $\geq$1,500--2,000 mIU/mL [discriminatory zone]), and if ectopic pregnancy is suspected: progesterone level (<5 ng/mL suggests non-viable pregnancy), dilation and curettage (D\&C) with pathologic examination for products of conception, and laparoscopy (gold standard for ectopic pregnancy diagnosis).
\section*{References}
\bibliographystyle{unsrtnat}
\bibliography{references}

\clearpage
